% =============================================================================
\subsection{Exigences d'utilisabilité}

\begin{exigence}{TPB-UTI-001}{Lecture d'un coup d'œil}[
    justification = {On consulte sa montre en passant, souvent en une ou deux secondes.},
    priorite      = {Essentielle},
    verification  = {D},
    phase         = {Transverse},
  ]
  L'application doit permettre de lire, sans autre action que son ouverture, la
  température de chaque pièce et celles qui chauffent.
\end{exigence}

\begin{exigence}{TPB-UTI-002}{Peu d'appuis pour agir}[
    justification = {Le geste le plus courant est de rejouer \enquote{21° pendant 1 h} dans
                     une pièce.},
    priorite      = {Essentielle},
    verification  = {D},
    phase         = {Transverse},
    note          = {Chemin visé~: ouvrir la pièce, ouvrir le pilotage, choisir le dernier
                     réglage. Le nombre maximal d'appuis reste à fixer.},
  ]
  L'application doit permettre de rejouer le dernier forçage d'une pièce en quelques
  appuis, depuis la liste des pièces.
\end{exigence}

\begin{exigence}{TPB-UTI-003}{Aucune saisie de texte}[
    justification = {Une montre n'a pas de clavier.},
    priorite      = {Essentielle},
    verification  = {I},
    phase         = {Transverse},
  ]
  L'application ne doit demander à l'utilisateur la saisie d'aucun texte sur la montre.
\end{exigence}

\begin{exigence}{TPB-UTI-004}{Usage à une main}[
    justification = {L'autre main porte le panier de linge ou la tasse.},
    priorite      = {Importante},
    verification  = {D},
    phase         = {Transverse},
  ]
  L'application doit pouvoir être utilisée entièrement avec les boutons de la montre,
  portée au poignet, d'une seule main.
\end{exigence}

\begin{exigence}{TPB-UTI-005}{Attente signalée}[
    justification = {Un écran qui ne réagit pas fait croire que le bouton n'a pas
                     marché.},
    priorite      = {Importante},
    verification  = {D},
    phase         = {Transverse},
  ]
  Pendant qu'elle attend la réponse de Tado, l'application doit montrer qu'elle est en
  cours d'échange.
\end{exigence}

\begin{exigence}{TPB-UTI-006}{Vocabulaire de Tado}[
    justification = {L'utilisateur connaît déjà Tado~; il ne doit pas apprendre d'autres
                     mots.},
    priorite      = {Souhaitable},
    verification  = {I},
    phase         = {Transverse},
  ]
  L'application devrait nommer les pièces, les modes et les états comme l'application
  mobile Tado.
\end{exigence}

% =============================================================================
\subsection{Exigences d'interface}
\label{sec:interfaces}

L'application échange avec trois interlocuteurs. Le tableau ci-dessous dit, pour chaque
échange, les exigences qui portent ce que l'application en attend.

\begin{table}[H]
  \centering\small
  \begin{tabularx}{\linewidth}{@{}l>{\raggedright\arraybackslash}p{32mm}X>{\raggedright\arraybackslash}p{38mm}@{}}
    \toprule
    \textbf{Réf.} & \textbf{Partie} & \textbf{Échange} & \textbf{Exigences} \\
    \midrule
    IF-01 & Compte Tado & L'application lit la situation de la maison (pièces, températures, humidité, chauffe, présence, alertes, historique) et la modifie (consignes, modes, présence). & \refexig{TPB-CPT-001} à \refexig{TPB-CPT-005}, \refexig{TPB-ACT-001} à \refexig{TPB-ACT-007} \\
    IF-02 & Montre de l'utilisateur & L'utilisateur lit l'écran, appuie sur les quatre boutons, reçoit une vibration lors d'une alerte. & \refexig{TPB-NAV-001} à \refexig{TPB-NAV-005}, \refexig{TPB-PLT-001} à \refexig{TPB-PLT-007}, \refexig{TPB-ALE-006} \\
    IF-03 & Téléphone de l'utilisateur & La montre n'atteint Tado que lorsqu'elle est à portée du téléphone auquel elle est associée. & \refexig{TPB-ACT-006}, \refexig{TPB-PIL-011} \\
    \bottomrule
  \end{tabularx}
  \caption{Interfaces de l'application}
  \label{tab:interfaces}
\end{table}

% =============================================================================
\subsection{Exigences de droits d'accès}

\begin{exigence}{TPB-SEC-001}{Accès limité au compte relié}[
    justification = {La maison est un lieu privé~; la montre ne doit montrer et commander
                     que celle de son propriétaire.},
    priorite      = {Essentielle},
    verification  = {T},
    phase         = {Transverse},
  ]
  L'application ne doit donner accès qu'à la maison du compte Tado auquel elle est reliée.
\end{exigence}

\begin{exigence}{TPB-SEC-002}{Pas de mot de passe sur la montre}[
    justification = {Une montre se perd ou se prête plus facilement qu'un téléphone.},
    priorite      = {Importante},
    verification  = {I},
    phase         = {Transverse},
  ]
  L'application ne doit ni demander ni afficher le mot de passe Tado de l'utilisateur sur la
  montre.
\end{exigence}

% =============================================================================
\subsection{Exigences de gestion de l'information}

\begin{exigence}{TPB-INF-001}{Dernière situation conservée}[
    justification = {Afficher quelque chose dès l'ouverture, même sans réseau.},
    priorite      = {Importante},
    verification  = {T},
    phase         = {Transverse},
  ]
  L'application doit conserver la dernière situation connue de la maison, d'une ouverture
  à l'autre.
\end{exigence}

\begin{exigence}{TPB-INF-002}{Dernier forçage conservé}[
    justification = {Rejouer le dernier réglage de chaque pièce (\refexig{TPB-PIL-004}).},
    priorite      = {Importante},
    verification  = {T},
    phase         = {Transverse},
  ]
  L'application doit conserver, pour chaque pièce, le dernier forçage demandé par
  l'utilisateur.
\end{exigence}

\begin{exigence}{TPB-INF-003}{Données limitées à l'usage}[
    justification = {Les habitudes d'une maison (heures de chauffe, présence) en disent long
                     sur ceux qui y vivent.},
    priorite      = {Importante},
    verification  = {I},
    phase         = {Transverse},
  ]
  L'application ne doit transmettre les données de la maison à personne d'autre que Tado
  et la montre de l'utilisateur, et ne doit conserver que ce qu'exigent ses écrans.
\end{exigence}
